% Generated by Listas/make_statements.py from pratica-lecture-05.md, item 18. Do not edit.
\textbf{Source.} Tadelis (2013), \emph{Game Theory: An Introduction}, Exercise 10.5 ``Negative Ad Campaigns Revisited'', pp.~216--217 {[}PDF 235--236{]}. Answer in the source: none (the book prints no solutions; the preface says about half are on Princeton UP's website, not checked). Reproduced verbatim.



\begin{quote}

\textbf{Negative Ad Campaigns Revisited:} Recall exercise 5.14 from Chapter 5, in which each one of two political parties can choose to buy time on commercial radio shows to broadcast negative ad campaigns against its rival. These choices are made simultaneously. Government regulations forbid a party from buying more than 2 hours of negative campaign time, so that each party cannot choose an amount of negative campaigning above 2 hours. Given a pair of choices \((a_1, a_2)\), the payoff of party \(i\) is given by the following function: \(v_i(a_1, a_2) = a_i - 2a_j + a_i a_j - (a_i)^2\).



\begin{enumerate}

\def\labelenumi{\alph{enumi}.}

\item

  Find the unique pure-strategy Nash equilibrium of the one-shot game.

\item

  If the parties could sign a binding agreement on how much to campaign, what levels would they choose?

\item

  Now assume that this game is repeated infinitely often, and the foregoing demonstrates the choices and payoffs per period. For which discount factors \(\delta \in (0, 1)\) can the levels you found in (b) be supported as a subgame-perfect equilibrium of the infinitely repeated game?

\item

  Despite the parties' ability to coordinate, as you have demonstrated in your answer to (c), the government is concerned about the parties' ability to run up to 2 hours a day of negative campaigning, and it is considering limiting negative campaigning to \(\frac{1}{2}\) hour a day, so that now \(a_i \in [0, \frac{1}{2}]\). Is this a good policy to limit negative campaigns further? Justify your answer with the relevant calculations. What is the intuition for your conclusion?

\end{enumerate}

\end{quote}

